% =====================================================================
%  4. Metodologia — como o trabalho foi feito
% =====================================================================
\section{Metodologia}
\label{sec:metodologia}

\indent\indent O procedimento adotado parte de uma frase correta, modificada de forma controlada. Cada 
frase variante introduz um único desvio, de tipo conhecido. Todos os métodos são aplicados às
mesmas frases, e observa-se o que cada um detecta, o que deixa passar e por qual mecanismo.

A frase de referência é \emph{Há três dias o incêndio atingiu uma área de mata nativa},
redigida a partir dos campos de um registro de eventos de fogo. Dela derivam as doze
variantes apresentadas na Tabela~\ref{tab:perturbacoes}.

\begin{table}[H]
\centering
\dimtable{
\begin{tabular}{@{}llp{5.4cm}@{}}
\toprule
\textbf{Variante} & \textbf{Nível afetado} & \textbf{Alteração introduzida} \\
\midrule
\texttt{natvia}                & ortografia   & forma inexistente no léxico \\
\texttt{A três dias}           & escolha lexical & forma existente, mas inadequada \\
\texttt{uma áreas}             & concordância nominal & determinante × núcleo \\
\texttt{a área inteiro}        & concordância nominal & adjetivo posposto \\
\texttt{os incêndios atingiu}  & concordância verbal & sujeito nominal × verbo \\
\texttt{cantou}                & semântica    & restrição de seleção violada \\
\texttt{que atingiu}           & estrutura    & relativa sem oração principal \\
\texttt{Atingir\ldots}         & estrutura    & infinitivo isolado \\
ordem embaralhada              & ordem        & as onze palavras permutadas \\
\texttt{tres} + \texttt{áreas} + \texttt{natvia} & múltiplo & três desvios simultâneos \\
\midrule
\texttt{destruiu}              & ---          & paráfrase plausível (controle) \\
\texttt{O incêndio é intenso}  & ---          & predicação nominal (controle) \\
\bottomrule
\end{tabular}}
\caption{As doze variantes da frase de referência.}
\label{tab:perturbacoes}
\end{table}

As duas últimas linhas não contêm desvio e destinam-se a verificar o inverso: um método que
as acusa produz falso positivo. A forma \texttt{destruiu} constitui paráfrase igualmente
correta, e \emph{O incêndio é intenso} é predicação nominal legítima, em que o verbo não
ocupa a raiz da árvore.

Quatro frases auxiliares completam o corpus. Duas comparam comprimento de dependência com as
mesmas palavras em ordens distintas, e duas formam um par de permutação. São necessárias
porque os descritores só se distinguem quando o comprimento é mantido fixo.

Os métodos distribuem-se por quatro dimensões, cada uma associada a uma verificação
distinta. A ortografia verifica se as palavras existem; a gramática, se elas concordam entre
si; a estrutura, se a frase dispõe de predicado e como suas partes se articulam; e a
semântica, se o enunciado é plausível. A divisão não é apenas organizativa, uma vez que cada
dimensão pressupõe a anterior: não se avalia concordância sobre uma palavra que não existe,
nem plausibilidade sobre uma frase desprovida de predicado. É essa dependência que determina
a ordem das etapas da \emph{pipeline}.

Nenhum modelo foi treinado ou ajustado. Os limites observados adiante pertencem a eles, e não ao método: a substituição do analisador sintático
altera o resultado sem alterar uma linha da regra.